\chapter{Communication}
In this chapter the communication between the customer (SeeSat e.V.) and Supplier (TSL23 TM Group) is described. The methods of information exchange are stated and cyclic intervals defined.

\section{Reporting System}
To inform the customer a reporting system via E Mail is established. The responsible on the supplier side writes a weekly E Mail containing information on project status with additional questions from the team. The questions part includes old not yet answered and new ones. The report includes a brief update on current and upcoming tasks.

\section{Action Monitoring}
The action monitoring is split into two categories. The first consists of monitoring general tasks and the second is the management of code writing. General tasks are tracked with action items (AI) who are conducted after each meeting. An AI is used to communicate the tasks each team member does at the moment. Therefore each AI is assigned a person. After working on a task the work is reviewed by other team members internally before the AIs deadline. Code related work items are monitored on GitLab using a ticket system further explained in the development process document.

\section{Customer-Supplier Meetings}
There are no regular meetings planed with the customer but there are regular update E Mails send by the supplier.

\section{Meeting Documentation}
In each meeting one person is responsible to write nodes that document the most important content discussed and to save it for later use. The responsible is chosen before the start of every meeting. The document will be saved to the corresponding meetings folder on the SeeSat Nextcloud in the meetings folder followed by the date the meeting was held.

\section{Responsibility}
In the weekly team meeting the content for each update mail is conducted. Afterwards additions like new risen questions and not yet answered questions (chapter Reporting System) are written by the communication responsible Nico. The E Mail is shared with the team members after its first iteration and additions or changes can be added to form a second iteration. After a “yes” from every team member the E Mail is send to the customer contacts Marcel and Stefan on the end of the week. Responses from the customer are shared as soon as possible with the team.

\section{Internal Team Communication}
The primary online communication tool used is a private discord channel. All event times are written there to be accepted or declined by every team member to centrally collect and inform every person. Status updates and requests for reviews are written in the chat. This allows every member to see the current problems or help each other.

\section{Document Exchange}
In the project lifetime many documents will be created. To exchange documents with the customer for explicit review a folder structure on the SeeSat Nextcloud created by the customer. Therefore a folder named "From SeeSat" is used the receive documents an a folder named "To SeeSat" to send documents. It is expected that both sides will read its corresponding received documents and communicate any changes that is wished to be made. Otherwise the content will be understood as acknowledged.